% \subsubsection{Absence de collecte de données en temps réel}
% \begin{itemize}
%     \item Actuellement, \textbf{aucune donnée n’est enregistrée ou stockée} automatiquement depuis les machines en salle de réveil.
%     \item La surveillance repose uniquement sur \textbf{un écran centralisé} situé dans une salle administrative.
%     \item \textbf{Les soignants doivent vérifier manuellement} l’état des patients à intervalles réguliers, ce qui peut entraîner des \textbf{retards} dans l’identification des urgences.
% \end{itemize}

\subsubsection{Manque de réactivité et d’optimisation des alertes}
\begin{itemize}
    \item Il n’existe \textbf{aucun système de gestion des alertes} priorisant les urgences en fonction de la gravité des paramètres vitaux.
    \item Les soignants \textbf{ne reçoivent pas d’alertes sur leurs appareils mobiles}, ce qui oblige à une surveillance constante de l’écran centralisé.
    \item L’absence d’un mécanisme de tri et de priorisation des alertes peut \textbf{retarder l’intervention} sur les cas critiques.
\end{itemize}

\subsubsection{Charge de travail accrue et inefficacité opérationnelle}
\begin{itemize}
    \item La \textbf{surveillance manuelle} et les déplacements constants du personnel entraînent une \textbf{perte de temps} et augmentent la charge de travail des soignants.
    \item \textbf{Risque d’erreurs humaines} dans la détection des anomalies, surtout en cas de forte affluence ou de fatigue du personnel.
    \item L'absence de \textbf{suivi automatisé des tendances des signes vitaux} empêche une anticipation des détériorations de l’état du patient.
\end{itemize}